\documentclass[10pt,aspectratio=169]{beamer}
\usepackage{bma01}

\title[Supremo, Bolzano y Weierstrass]{Ínfimo y supremo de un subconjunto de $\R$}
\subtitle{Teoremas de continuidad: valor intermedio de Bolzano y de Weierstrass\\[2pt]\small Semana 5 -- Sesión 2}

\begin{document}

\begin{frame}
  \titlepage
\end{frame}

\begin{frame}{Objetivos de la sesión}
  \begin{enumerate}
    \item Definir cotas, \resaltar{supremo} e \resaltar{ínfimo}; distinguirlos de máximo y mínimo.
    \item Comprender el \resaltar{axioma del supremo} (completitud de $\R$).
    \item Enunciar y aplicar el \resaltar{teorema de Bolzano} y el \resaltar{teorema del valor intermedio}; método de bisección.
    \item Enunciar y aplicar el \resaltar{teorema de Weierstrass} (valores extremos).
  \end{enumerate}
  \vspace{1mm}
  {\small\textbf{Competencia 2 del sílabo:} Evalúa propiedades de límites y continuidad y su utilidad en aplicaciones de ciencia e ingeniería.}
  \begin{block}{Plan (2 horas; incluye verificación previa de 5 min)}
    \begin{tabular}{ll}
      0:00 -- 0:35 & Cotas, supremo e ínfimo. Axioma del supremo\\
      0:35 -- 1:10 & Teorema de Bolzano y del valor intermedio. Bisección\\
      1:10 -- 1:40 & Teorema de Weierstrass. Importancia de las hipótesis\\
      1:40 -- 2:00 & Actividad colaborativa, resumen y cierre metacognitivo
    \end{tabular}
  \end{block}
\end{frame}

\begin{frame}{Aula invertida: preparación y verificación}
  \small
  \begin{previa}[antes de la sesión]
    \begin{enumerate}
      \item Lea sobre cotas, supremo e ínfimo y los teoremas de Bolzano y Weierstrass.
      \item Halle $\sup$ e $\inf$ de $\{1/n:n\in\N\}$ y de $(0,1)$.
      \item Pregunta guía: ¿qué falla en Bolzano si la función no es continua?
    \end{enumerate}
  \end{previa}
  \begin{block}{Verificación al inicio (5 min, en pares)}
    \begin{enumerate}
      \item ¿Tiene máximo el conjunto $(0,1)$? ¿Y supremo?
      \item Enuncie las hipótesis del teorema de Bolzano.
    \end{enumerate}
    Los pares discuten y votan; el docente aclara solo lo que genera dudas y ajusta el ritmo de la sesión.
  \end{block}
\end{frame}

\section{Supremo e ínfimo}

\begin{frame}{Cotas, supremo e ínfimo}
  Sea $A\subseteq\R$, $A\neq\emptyset$.
  \begin{definicion}
    \begin{itemize}
      \item $M$ es \resaltar{cota superior} de $A$ si $a\le M$ para todo $a\in A$. $A$ es \emph{acotado superiormente} si tiene una.
      \item $s=\sup A$ (\resaltar{supremo}) es la \textbf{menor} cota superior: (i) $s$ es cota superior y (ii) si $M$ es cota superior, $s\le M$.
      \item Análogamente: cota inferior e \resaltar{ínfimo} $\inf A$ (la mayor cota inferior).
      \item Si $\sup A\in A$, se llama \resaltar{máximo}; si $\inf A\in A$, \resaltar{mínimo}.
    \end{itemize}
  \end{definicion}
  \begin{center}
    \begin{tikzpicture}[scale=1.1]
      \draw[->] (-0.5,0) -- (9.5,0);
      \draw[uniAzul, line width=3pt] (1.5,0) -- (5,0);
      \draw[uniAzul, thick, fill=white] (1.5,0) circle (2.5pt);
      \fill[uniAzul] (5,0) circle (2.5pt);
      \node[above, uniAzul] at (3.2,0.1) {\scriptsize $A=(1,4]$};
      \node[below] at (1.5,-0.15) {\scriptsize $\inf A=1$ (no es mínimo)};
      \node[below] at (5,-0.15) {\scriptsize $\sup A=\max A=4$};
      \foreach \x in {6,7,8.5} \draw[uniGranate, thick] (\x,0.12) -- (\x,-0.12);
      \node[above, uniGranate] at (7.2,0.1) {\scriptsize cotas superiores};
    \end{tikzpicture}
  \end{center}
\end{frame}

\begin{frame}{Caracterización y ejemplos}
  \begin{teorema}[caracterización con $\varepsilon$]
    Si $s$ es cota superior de $A$: \quad $s=\sup A\iff\forall\varepsilon>0\ \exists a\in A:\ a>s-\varepsilon$.
  \end{teorema}
  \small (Ningún número menor que $s$ es cota superior: siempre hay elementos de $A$ ``muy cerca'' de $s$.)
  \normalsize
  \begin{ejemplo}
    \begin{enumerate}
      \item $A=\left\{1-\frac1n:n\in\N\right\}=\{0,\frac12,\frac23,\dots\}$: \ $\inf A=\min A=0$, \ $\sup A=1\notin A$.
        \ Dado $\varepsilon$, basta $n>\frac1\varepsilon$.
      \item $A=\{x\in\R:x^2<2\}=(-\sqrt2,\sqrt2)$: \ $\sup A=\sqrt2$, sin máximo.
      \item $\N$ no es acotado superiormente (propiedad arquimediana).
      \item $A=\{\sen x:x\in\R\}$: \ $\sup A=\max A=1$.
    \end{enumerate}
  \end{ejemplo}
\end{frame}

\begin{frame}{El axioma del supremo}
  \begin{teorema}[Axioma del supremo (completitud)]
    Todo subconjunto no vacío de $\R$ acotado superiormente tiene supremo \textbf{en $\R$}.
  \end{teorema}
  \begin{ejemplo}[$\Q$ no es completo]
    $B=\{q\in\Q:\ q>0,\ q^2<2\}$ está acotado en $\Q$, pero no tiene supremo en $\Q$: su supremo en $\R$ es $\sqrt2\notin\Q$.
  \end{ejemplo}
  \begin{idea}
    La completitud es lo que distingue a $\R$ de $\Q$: la recta real ``no tiene huecos''. Es la base de los teoremas de esta sesión
    (y de casi todo el cálculo).
  \end{idea}
  \begin{observacion}
    Para funciones: $\sup_{x\in I}f(x)=\sup\{f(x):x\in I\}$. \ Ej.: $\sup_{x>0}\frac{x}{x+1}=1$ (no se alcanza).
  \end{observacion}
\end{frame}

\section{Teorema de Bolzano y del valor intermedio}

\begin{frame}{Teorema de Bolzano}
  \begin{columns}
    \begin{column}{0.55\textwidth}
      \begin{teorema}[Bolzano]
        Si $f$ es continua en $[a,b]$ y $f(a)\cdot f(b)<0$, entonces existe $c\in(a,b)$ con $f(c)=0$.
      \end{teorema}
      \textbf{Idea de la prueba} ($f(a)<0<f(b)$). Sea $A=\{x\in[a,b]:f(x)<0\}$, no vacío y acotado; $c=\sup A$.
      \begin{itemize}
        \item Si $f(c)<0$, por continuidad $f<0$ un poco a la derecha de $c$: contradice que $c$ es cota superior.
        \item Si $f(c)>0$, $f>0$ cerca de $c$: entonces $c-\delta$ sería cota menor. Contradicción.
      \end{itemize}
      Luego $f(c)=0$. \ $\blacksquare$
    \end{column}
    \begin{column}{0.45\textwidth}
      \centering
      \begin{tikzpicture}
        \begin{axis}[bma, xmin=0, xmax=4, ymin=-3, ymax=2.5, width=6.5cm, height=5.5cm, xtick={0.5,2.17,3.5}, xticklabels={$a$,$c$,$b$}, ytick=\empty, samples=100, domain=0.5:3.5]
          \addplot[curva] {0.4*(x-2.17)^3+0.3*(x-2.17)+0.2*sin(deg(4*x))*(x-2.17)};
          \addplot[only marks, mark=*, uniGranate] coordinates {(2.17,0)};
          \addplot[only marks, mark=*, uniAzul] coordinates {(0.5,-2.668) (3.5,1.604)};
        \end{axis}
      \end{tikzpicture}

      \small Una curva continua que pasa de negativa a positiva debe cruzar el eje.
    \end{column}
  \end{columns}
\end{frame}

\begin{frame}{Método de bisección}
  \begin{columns}
    \begin{column}{0.55\textwidth}
      Bolzano sugiere un algoritmo: partir el intervalo por la mitad y quedarse con la mitad donde hay cambio de signo.
      \begin{ejemplo}[1]
        $f(x)=x^3-x-1$ en $[1,2]$: $f(1)=-1<0$, $f(2)=5>0$.
        \begin{center}
          \small
          \begin{tabular}{c|c|c|c}
            $a$ & $b$ & $m$ & $f(m)$\\ \hline
            1 & 2 & 1.5 & $+0.875$\\
            1 & 1.5 & 1.25 & $-0.297$\\
            1.25 & 1.5 & 1.375 & $+0.225$\\
            1.25 & 1.375 & 1.3125 & $-0.052$\\
            1.3125 & 1.375 & 1.34375 & $+0.083$\\
            1.3125 & 1.34375 & 1.328125 & $+0.015$
          \end{tabular}
        \end{center}
        Raíz $\approx1.3247$. Tras $n$ pasos, el error es $\le\frac{b-a}{2^n}$.
      \end{ejemplo}
    \end{column}
    \begin{column}{0.45\textwidth}
      \centering
      \begin{tikzpicture}
        \begin{axis}[bma, xmin=0.9, xmax=2.1, ymin=-1.5, ymax=5.5, width=6.3cm, height=5.6cm, xtick={1,1.25,1.5,2}, ytick=\empty, samples=100, domain=0.95:2.05]
          \addplot[curva] {x^3-x-1};
          \addplot[only marks, mark=*, uniGranate] coordinates {(1.5,0.875) (1.25,-0.297) (1.375,0.225)};
          \addplot[only marks, mark=*, uniVerde] coordinates {(1.3247,0)};
        \end{axis}
      \end{tikzpicture}
    \end{column}
  \end{columns}
\end{frame}

\begin{frame}{Teorema del valor intermedio y aplicaciones}
  \begin{teorema}[Valor intermedio]
    Si $f$ es continua en $[a,b]$ y $k$ está entre $f(a)$ y $f(b)$, existe $c\in[a,b]$ con $f(c)=k$.
  \end{teorema}
  \textbf{Prueba.} Aplique Bolzano a $g(x)=f(x)-k$. \ $\blacksquare$
  \begin{ejemplo}[2: polinomios de grado impar]
    Todo polinomio de grado impar tiene al menos una raíz real: $P(x)\to\pm\infty$ con signos opuestos en $\pm\infty$.
  \end{ejemplo}
  \begin{ejemplo}[3: punto fijo]
    $\cos x=x$ tiene solución en $[0,1]$: $g(x)=\cos x-x$, $g(0)=1>0$, $g(1)=\cos1-1<0$. \ ($x^*\approx0.7391$.)
  \end{ejemplo}
  \begin{ejemplo}[4: física]
    Si la temperatura en un anillo varía continuamente, existen dos puntos diametralmente opuestos con la misma temperatura:
    $g(\theta)=T(\theta)-T(\theta+\pi)$ cumple $g(\pi)=-g(0)$.
  \end{ejemplo}
\end{frame}

\begin{frame}{¡Las hipótesis importan!}
  \begin{columns}
    \begin{column}{0.5\textwidth}
      \begin{ejemplo}[sin continuidad]
        $f(x)=\sgn(x)$ en $[-1,1]$: $f(-1)=-1$, $f(1)=1$, pero nunca vale $\frac12$.
      \end{ejemplo}
      \begin{ejemplo}
        $\frac1x$ en $[-1,1]$: cambia de signo y no se anula; no es continua en $0$.
      \end{ejemplo}
      \begin{observacion}
        El recíproco es falso: $f$ puede tener la propiedad del valor intermedio sin ser continua (p.\,ej.\ $\sen\frac1x$ con $f(0)=0$).
      \end{observacion}
    \end{column}
    \begin{column}{0.5\textwidth}
      \centering
      \begin{tikzpicture}
        \begin{axis}[bma, xmin=-1.3, xmax=1.3, ymin=-1.5, ymax=1.5, width=6.3cm, height=5cm, xtick={-1,1}, ytick={-1,0.5,1}]
          \addplot[curva] coordinates {(-1,-1) (0,-1)};
          \addplot[curva] coordinates {(0,1) (1,1)};
          \draw[uniGranate, dashed] (axis cs:-1.3,0.5) -- (axis cs:1.3,0.5);
          \node[uniGranate] at (axis cs:-0.7,0.7) {\scriptsize $y=\frac12$ nunca se alcanza};
        \end{axis}
      \end{tikzpicture}
    \end{column}
  \end{columns}
\end{frame}

\section{Teorema de Weierstrass}

\begin{frame}{Teorema de Weierstrass (valores extremos)}
  \begin{columns}
    \begin{column}{0.55\textwidth}
      \begin{teorema}[Weierstrass]
        Si $f$ es continua en un intervalo \textbf{cerrado y acotado} $[a,b]$, entonces $f$ es acotada y alcanza un máximo y un mínimo:
        existen $x_1,x_2\in[a,b]$ con
        \[ f(x_1)\le f(x)\le f(x_2)\quad\forall x\in[a,b]. \]
      \end{teorema}
      \begin{proposicion}[consecuencia]
        La imagen de $[a,b]$ por una función continua es un intervalo cerrado $[m,M]$ (por Weierstrass y el TVI).
      \end{proposicion}
      \small En la Semana 9 usaremos este teorema para garantizar que un problema de optimización tiene solución.
    \end{column}
    \begin{column}{0.45\textwidth}
      \centering
      \begin{tikzpicture}
        \begin{axis}[bma, xmin=0, xmax=4.3, ymin=0, ymax=3.3, width=6.5cm, height=5.5cm, xtick={0.4,1.35,3.42,4}, xticklabels={$a$,$x_2$,$x_1$,$b$}, ytick=\empty, samples=100, domain=0.4:4]
          \addplot[curva] {1.7+0.9*sin(deg(1.4*x-0.2))+0.15*x};
          \addplot[only marks, mark=*, uniGranate] coordinates {(1.35,2.80) (3.42,1.32)};
          \draw[dashed, uniGranate] (axis cs:0,2.80) -- (axis cs:4.3,2.80);
          \draw[dashed, uniGranate] (axis cs:0,1.32) -- (axis cs:4.3,1.32);
          \node[uniGranate] at (axis cs:2.3,3.05) {\scriptsize máximo $M$};
          \node[uniGranate] at (axis cs:2.3,1.1) {\scriptsize mínimo $m$};
        \end{axis}
      \end{tikzpicture}
    \end{column}
  \end{columns}
\end{frame}

\begin{frame}{Las hipótesis de Weierstrass son necesarias}
  \begin{columns}
    \begin{column}{0.33\textwidth}
      \centering
      \begin{tikzpicture}
        \begin{axis}[bma, xmin=0, xmax=1.2, ymin=0, ymax=8, width=4.8cm, height=4.3cm, xtick={1}, ytick=\empty, samples=100]
          \addplot[curva, domain=0.13:1] {1/x};
          \addplot[only marks, mark=*, uniAzul] coordinates {(1,1)};
        \end{axis}
      \end{tikzpicture}\\
      \small $\frac1x$ en $(0,1]$: no es cerrado; no acotada.
    \end{column}
    \begin{column}{0.33\textwidth}
      \centering
      \begin{tikzpicture}
        \begin{axis}[bma, xmin=-0.1, xmax=1.2, ymin=0, ymax=1.3, width=4.8cm, height=4.3cm, xtick={1}, ytick={1}]
          \addplot[curva, domain=0:1] {x};
          \addplot[only marks, mark=*, uniAzul] coordinates {(0,0)};
          \addplot[only marks, mark=o, uniAzul, fill=white] coordinates {(1,1)};
        \end{axis}
      \end{tikzpicture}\\
      \small $x$ en $[0,1)$: $\sup=1$ no se alcanza.
    \end{column}
    \begin{column}{0.33\textwidth}
      \centering
      \begin{tikzpicture}
        \begin{axis}[bma, xmin=-0.1, xmax=1.2, ymin=0, ymax=1.3, width=4.8cm, height=4.3cm, xtick={1}, ytick={1}]
          \addplot[curva, domain=0:0.98] {x};
          \addplot[only marks, mark=o, uniAzul, fill=white] coordinates {(1,1)};
          \addplot[only marks, mark=*, uniAzul] coordinates {(1,0)};
        \end{axis}
      \end{tikzpicture}\\
      \small $f(x)=x$, $f(1)=0$ en $[0,1]$: discontinua; sin máximo.
    \end{column}
  \end{columns}
  \begin{ejemplo}[5]
    $f(x)=x^3-3x$ en $[-2,3]$: continua en un cerrado, así que tiene máximo y mínimo absolutos. (Los hallaremos con derivadas:
    $\max=f(3)=18$, $\min=f(-2)=f(1)=-2$.)
  \end{ejemplo}
\end{frame}

\begin{frame}{Actividad colaborativa: Bisección en Python}
  \small
  \begin{activa}[Aprendizaje basado en la experiencia]
    Demuestre que $\cos x=x$ tiene una solución en $[0,1]$. Luego \textbf{programe} (Python u hoja de cálculo) el método de bisección, haga 10 iteraciones, estime la raíz y acote el error con $\frac{b-a}{2^n}$. Compare con la solución de un CAS.
  \end{activa}
  \begin{itemize}
    \item \textbf{Organización:} equipos de 3 a 4 con roles rotativos (coordinador, redactor, verificador, portavoz); 12--15 minutos.
    \item \textbf{Producto:} Código, tabla de iteraciones y justificación de la cota.
    \item \textbf{Retroalimentación:} el docente recorre los equipos, hace preguntas y no da la solución hasta la puesta en común.
  \end{itemize}
\end{frame}

\section{Ejercicios}

\begin{frame}{Ejercicios propuestos}
  \small
  \begin{ejercicio}
    \begin{enumerate}
      \item Halle supremo, ínfimo, máximo y mínimo (si existen): (a) $\left\{\frac{n}{n+1}:n\in\N\right\}$, \ (b) $\left\{\frac{(-1)^n}{n}:n\in\N\right\}$,
        \ (c) $\{x:x^2-3x<4\}$, \ (d) $\left\{\frac{x^2}{1+x^2}:x\in\R\right\}$.
      \item Pruebe que $\sup\left\{1-\frac1{2^n}\right\}=1$ usando la caracterización con $\varepsilon$.
      \item Pruebe que si $A\subseteq B$ y $B$ es acotado, entonces $\sup A\le\sup B$.
      \item Pruebe que $x^5-3x+1=0$ tiene al menos tres raíces reales.
      \item Aproxime con 4 pasos de bisección la raíz de $e^{-x}=x$ en $[0,1]$.
      \item Pruebe que si $f:[0,1]\to[0,1]$ es continua, existe $c$ con $f(c)=c$ (punto fijo).
      \item ¿Tiene máximo $f(x)=\frac{x}{1+x^2}$ en $[0,\infty)$? (Weierstrass no se aplica directamente; analice.)
      \item Un excursionista sube una montaña de 6:00 a 18:00 y baja por el mismo camino al día siguiente, de 6:00 a 18:00. Pruebe que hay
        un instante en que estuvo en el mismo lugar a la misma hora ambos días.
    \end{enumerate}
  \end{ejercicio}
\end{frame}

\begin{frame}{Metacognición y uso responsable de la IA}
  \small
  \begin{columns}[t]
    \begin{column}{0.5\textwidth}
      \begin{metacog}[cierre, 3 min]
        \begin{enumerate}
          \metacomun
          \item ¿Puedo explicar por qué la bisección converge sin invocar al código? Escríbalo.
        \end{enumerate}
      \end{metacog}
    \end{column}
    \begin{column}{0.5\textwidth}
      \begin{iabox}
        Pida a una IA el código de bisección. Léalo línea por línea, pruébelo con una función de raíz conocida e indique qué cambió.\\[2pt]
        \textit{Regla:} primero intente usted; luego use la herramienta para contrastar, nunca para sustituir su razonamiento. Anote el cambio en su bitácora.
      \end{iabox}
    \end{column}
  \end{columns}
\end{frame}

\begin{frame}{Resumen}
  \begin{itemize}
    \item $\sup A$: menor cota superior; si pertenece a $A$, es el máximo. Axioma del supremo: $\R$ es completo.
    \item \resaltar{Bolzano}: continua en $[a,b]$ y cambio de signo $\Rightarrow$ hay una raíz. Bisección: error $\le\frac{b-a}{2^n}$.
    \item \resaltar{TVI}: una función continua toma todos los valores intermedios.
    \item \resaltar{Weierstrass}: continua en $[a,b]$ $\Rightarrow$ alcanza máximo y mínimo.
    \item Sin continuidad o sin intervalo cerrado y acotado, las conclusiones pueden fallar.
  \end{itemize}
  \vspace{3mm}
  \textbf{Próxima sesión:} la derivada; rectas tangente y normal.
\end{frame}


\begin{frame}{Evaluación y trabajo posterior}
  \small
  \begin{evalbox}[Práctica del miércoles de la semana 5: PD3]
    \textbf{Cobertura:} semanas 4 y 5. \textbf{Naturaleza:} práctica dirigida en equipo, con retroalimentación inmediata.
  \end{evalbox}
  \begin{itemize}
    \item \criterios
    \item \declaracionia
    \item \textbf{Portafolio:} guarde el producto de la actividad, la bitácora de metacognición y los ejercicios propuestos de esta sesión; son evidencia de progreso y base de la retroalimentación.
  \end{itemize}
\end{frame}

\end{document}
